% BEGIN D85G1
\paragraph{D85G1. Wild inertia of every original localized unitary label at a type prime.}
Fix an original member $D'$, a location $i$, a prime $w\in S_i$
above $g=g_k$, and any $P\in\mathrm{AV}(i)$. Put $n=r+t=n_V$.
Every prime above $g$ splits over the maximal real field; $w$ is
either a designated $w_0=\mathrm{plGw}(k,v)$ or its conjugate.
Both completions are $\mathbb Q_g$. We show that for every
$\sigma\in\mathrm{wildSet}(D'.\mathrm{nd},D'.\mathrm{nd}.\mathrm{plOm}_w)$
one has
\[
   \mathrm{inePol}\bigl(\mathrm{locV}(i,P,w),\sigma\bigr)=(X-1)^n .
\]

Use the member's level $M_i$ and the preferred split coordinate
$s$ for $w_0$. Its tuple index agrees with $\mu_{\rm Tw}$ of the
coefficient place by the localization identity, and its zero-based
row is $j(i)-1$. Thus the original input $\mathrm{mt}.\mathrm{uni}$
applies at this $k,v$. Write $K_1$ for the kernel of reduction
$\mathrm{GL}_n(\mathbb Z_g)\to\mathrm{GL}_n(\mathbb F_g)$.
Given $h\in K_1$, insert $h$ in the local split factor and the
identity at every other real finite place. Its other CM component
is the inverse transpose of $h$ in the paired frame, as follows
from the unitary equation. The resulting adele is supported at
the one tested place. In the preferred coordinate its lower-left
block is zero modulo $g$, and its upper-left block is the identity
modulo $g$. This gives precisely the inequalities for
$\mathrm{UniRT}$, so the adele belongs to the level.
The character there is the product of the additive character of the
sum of the upper-left superdiagonal entries and the multiplicative
character of the lower-right determinant. The sum lies in
$g\mathbb Z_g$, on which the specified additive character is
trivial. The determinant belongs to $1+g\mathbb Z_g$.
The character on this block is the local component of the tuple's
$\mu^0_j\xi$: both factors are tame by the member's character frame
and $\mathrm{tp}.\mathrm{xi\_tame}$. Indeed the local component
test agrees on the dense global multiplicative group and the two
characters are smooth; hence the determinant factor also equals
one. At the conjugate CM place, the completion isomorphism carries
integer rings and maximal ideals to each other, and inverse
transpose bijects $K_1$ with itself. The same single-place adeles
therefore realize the full $K_1$ in either split coordinate, with
trivial level character.

Choose a nonzero member of $\mathrm{typeV}(i,P)$, which is nonzero
for every $P$, including those without a Hecke point. The injective
map $M_i.\mathrm{real}$ sends it to a nonzero level-character form.
The member's $\mathrm{Au}.\mathrm{typeV}$ puts it in the isotypic
sum of its actual irreducible representation $U_P$. It is fixed by
the preceding $K_1$. Projecting to a copy of $U_P$ detecting it
gives a nonzero fixed vector there: an algebraic sum of simple
submodules is a direct sum of a suitable family of such modules.
At the chosen $w$, the restriction of $U_P$ is isotypic for its
local component $\pi_w$. Indeed the translates away from $w$ of an
embedded copy of $\pi_w$ span a nonzero globally stable subspace
and hence all of $U_P$. Project once more to a copy detecting the
vector. This proves $\pi_w^{K_1}\ne0$ for the original label $P$.

For completeness the support of $\pi_w$ is good in the precise
sense required by the weak correspondence. Choose its supercuspidal
support and realize it as a subquotient of the corresponding
normalized parabolic induction. Compact averaging in characteristic
zero is exact, so the induction has nonzero $K_1$-invariants.
By Iwasawa decomposition its restriction to the maximal compact
is induced from the integral parabolic. Since $K_1$ is normal
in the maximal compact, a nonzero value of a fixed function
belongs to the tensor product of the principal-congruence fixed
spaces of the Levi factors. Each support factor thus has such a
fixed vector. Over $\mathbb Q_g$ this is an allowed good support
factor, taking the identity character twist.
The depth law of the member's local correspondence now makes its
$n$-dimensional Weil--Deligne parameter trivial on wild inertia.
The original comparison $\mathrm{Au}.\mathrm{emb}$, the native
$\mathrm{embAV\_loc}$ identity and the split local realization
identify this parameter with the image of $\mathrm{locV}(i,P,w)$
under $A.\mathrm{wdT}$, as in B85U1. For $\sigma$ in the stated
wild set, the condition $|\sigma(y)/y-1|<1$ for nonzero $y$ is
equivalently $|\sigma(y)-y|<|y|$, and puts $\sigma$ in inertia
with the required wild condition. The inertia clause of
$A.\mathrm{wdT}$ thus reads the characteristic polynomial of the
identity in dimension $n$. This proves the original
$\mathrm{GaloisPreconditions.depthG}$ for every tested location.
% END D85G1
